\documentclass[11pt]{article}

\usepackage[margin=0.85in]{geometry}
\usepackage{fontspec}
\usepackage{microtype}
\usepackage{booktabs}
\usepackage{graphicx}
\usepackage{xcolor}
\usepackage{hyperref}
\usepackage{enumitem}

\setmainfont{Georgia}
\setsansfont{Segoe UI}
\setmonofont{Consolas}

\definecolor{ink}{HTML}{1F2937}
\definecolor{accent}{HTML}{0F766E}
\definecolor{accentlight}{HTML}{E6FFFB}

\hypersetup{colorlinks=true,linkcolor=accent,urlcolor=accent}
\pagestyle{empty}
\setlength{\parindent}{0pt}
\setlength{\parskip}{0.55em}

\begin{document}
\color{ink}

{\sffamily
\begin{center}
  {\Large\bfseries PushT Follow-Up Addendum}\\[0.3em]
  {\large AF-LeWM-lite v2 Sweep and Stability Check}\\[0.4em]
  {\small April 10, 2026}
\end{center}
}

\fcolorbox{accent}{accentlight}{%
  \parbox{\dimexpr\linewidth-2\fboxsep-2\fboxrule\relax}{
    \textbf{Updated conclusion.} The original AF-LeWM-lite v2 objective still gives the best 50-episode PushT control result, but the advantage over baseline is small and unstable. A second 50-episode evaluation seed reduces the gap to \textbf{5/100} successes for v2 versus \textbf{4/100} for baseline. The current evidence supports the existence of a real representation-control tradeoff, and it does not yet support a strong robustness win.
  }
}

\section*{Why this follow-up}
The first official PushT report showed a tension:
\begin{itemize}[leftmargin=1.2em, itemsep=0.2em, topsep=0.2em]
  \item baseline had the best shared JEPA objective,
  \item original v2 had the best 50-episode success,
  \item the gap was only one success: \texttt{3/50} versus \texttt{2/50}.
\end{itemize}

This addendum tests whether the adversarial nuisance term is responsible for both the control gain and the optimization damage.

\section*{Fair comparison metric}
Cross-model loss comparison uses the shared objective
\[
\mathcal{L}_{\mathrm{core}} = \mathcal{L}_{\mathrm{pred}} + 0.09\,\mathcal{L}_{\mathrm{sigreg}},
\]
because total training loss includes AF-specific auxiliary terms and is not directly comparable once those terms are active.

\section*{Sweep results}
\begin{center}
\small
\begin{tabular}{lrrl}
\toprule
Model & Core loss & Success (\%) & Comment \\
\midrule
Baseline LeWM & 0.1469 & 4.0 & reference \\
AF-LeWM-lite v2 & 0.4726 & 6.0 & original best control \\
v2 no adversary & 0.1491 & 0.0 & removed dynamics nuisance adversary \\
v2 soft & 0.3603 & 2.0 & reduced all auxiliary weights \\
v2 grl half & 0.3145 & 2.0 & halved GRL only \\
\bottomrule
\end{tabular}
\end{center}

\section*{What changed}
\textbf{Removing the adversary restores the shared world-model loss and destroys control.} The no-adversary run lands almost on top of baseline in core loss, then falls to \texttt{0/50} successes. That is the strongest sign that the dynamics-side nuisance adversary is doing real work for control.

\textbf{Weakening the adversary improves the loss and weakens control.} Both the soft variant and the GRL-half variant move the core loss back toward baseline, and both drop to \texttt{1/50} successes.

\textbf{The current tradeoff is sharp.} Better shared modeling and better PushT control are pulling in different directions under this short budget.

\section*{100-episode check}
I reran baseline and the original v2 on a second evaluation seed while keeping the same 50-episode protocol per seed.

\begin{center}
\small
\begin{tabular}{lrrr}
\toprule
Model & Seed 42 & Seed 43 & Aggregate \\
\midrule
Baseline LeWM & 2/50 & 2/50 & 4/100 = 4.0\% \\
AF-LeWM-lite v2 & 3/50 & 2/50 & 5/100 = 5.0\% \\
\bottomrule
\end{tabular}
\end{center}

This keeps v2 slightly ahead, and the margin is too small to claim a robust improvement.

\section*{Research conclusion}
\begin{itemize}[leftmargin=1.2em, itemsep=0.2em, topsep=0.2em]
  \item The dynamics nuisance adversary is the key mechanism behind the current v2 control gain.
  \item The same mechanism is also the main source of degradation in the shared JEPA objective.
  \item The present PushT evidence supports a \textbf{tradeoff}, not a clear overall win.
  \item The next experiment should tune \textbf{when} the adversary turns on, not only \textbf{how strong} it is. A delayed or ramped adversary schedule is the highest-value next test.
\end{itemize}

In one sentence: \textbf{the factorized-latent idea still looks alive, but the current v2 benefit is fragile and the right next move is schedule tuning rather than static weight reduction.}

\end{document}
